\BourbakiStarChapter{第二、三章历史注记}

（注：方括号中的数字指本注文末的参考文献。）

线性代数是数学中最古老又最新颖的分支之一。一方面，数学的起源涉及那些通过一次乘法或除法即可解决的问题，即通过计算某个函数 \(f(x) = ax\) 的值，或通过求解方程 ax = b：这些是线性代数的典型问题，若不“线性地思考”，就不可能处理它们，甚至无法正确地提出它们。

另一方面，不仅这些问题，而且几乎所有关于一次方程的内容，早已被归入基础教学之中，直到域、环、拓扑向量空间等概念的现代发展，才将线性代数中的本质概念（例如对偶性）分离并凸显出来；随后人们认识到，几乎整个现代数学本质上都具有线性特征，而“线性化”本身就是其显著特征之一，线性代数也因此获得了应有的地位。从我们当前的观点来撰写其历史，将是一项既困难又重要的任务；因此，我们只能满足于给出一个简要的概述。

由上可见，线性代数无疑是为了满足实际计算者的需要而诞生的；因此，我们看到三率法†和假位法，以或明或暗的表述，在从埃及的莱因德纸草书，经由阿耶波多、阿拉伯人、伦纳德，直到我们小学所用的所有实用算术手册中都扮演着重要角色。

比萨以及中世纪和文艺复兴时期无数本“计算手册”皆是如此；但这些内容从未构成最先进科学理论中超过一小部分的部分，仅供实际从业者使用。

至于数学家本身，他们对线性代数的研究性质取决于其科学的总体结构。古希腊数学，如欧几里得《几何原本》所阐述的，发展了两个具有线性特征的抽象理论，一方面是量的理论（{[}2{]}，第五卷；参见《一般拓扑学》第四卷的历史注记），另一方面是整数的理论（{[}2{]}，第七卷）。在巴比伦人那里，我们发现的方法更接近我们的初等代数；他们知道如何求解一次方程组，而且求解得非常优雅（{[}1{]}，第181—183页）。然而，在很长一段时间里，线性代数的进展主要局限于代数计算的范畴，应当从这个角度来考虑，这与本注记无关；要将一个线性方程组化为形如 ax = b 的方程，在单个未知量的情形下，只需知道将项从一边移到另一边以及合并同类项的规则（这些规则实质上已由丢番图陈述过）；而在多个未知量的情形下，只需知道如何依次消去它们，直到只剩下一个未知量。因此，直到十八世纪，代数学的论著都认为，只要阐述了这些规则，关于一次方程的一切就已完成；至于方程个数与未知量个数相同（他们不考虑其他情形）而左边不是线性无关形式的方程组，他们只是顺带指出这表明问题提得不恰当。在十九世纪的论著乃至某些更近期的著作中，这一观点仅仅因记号的进步而有所改变——记号使得能够写出 n 个未知量的 n 个方程组成的方程组——并且因行列式的引入而有所改变——行列式使得能够在“一般情形”下给出显式解的公式；这一进步，其功劳本应归于莱布尼茨（{[}7{]}，第239页），倘若他发展并发表了他关于这一主题的思想的话，但这一进步主要归功于十八世纪和十九世纪初的数学家们。

但我们必须首先研究各种思潮，它们对线性代数的发展所作的贡献远大于线性方程的研究，其意义正如我们今天所理解的那样。受阿波罗尼奥斯研究的启发，费马 \([4(a)]\) 在甚至先于笛卡尔 {[}5{]}构想了解析几何的原理之后，便有了按次数对平面曲线进行分类的想法（这一想法逐渐为所有数学家所熟悉，可以认为在十七世纪末已被明确掌握），并表述了基本原理：在平面中，一次方程表示一条直线，二次方程表示一条圆锥曲线；他由此立即推出一些关于几何轨迹的“非常优美”的推论。与此同时，他陈述了 \([4(b)]\) 将问题分为具有单一解的问题、归结为含两个未知量的方程的问题、含三个未知量的方程的问题等等；并且他补充道：第一类在于确定一个点，第二类在于确定一条直线或平面轨迹，其余各类在于确定一个曲面，等等（“\ldots{} 这样的问题不仅寻求一个点或一条线，而是寻求与问题相适应的整个曲面；在这里，曲面作为轨迹有其起源，其余情形亦如此”，同上引书，第186页；这里已经孕育了 n 维几何的萌芽）。这篇论文表述了代数与代数几何中的维数原理，表明代数与几何的融合完全符合现代观念，但正如已经看到的，这种融合花了两个多世纪才深入人心。

至少这些思想很快导致了解析几何的扩展，并在十八世纪随着克莱罗、欧拉、克拉默、拉格朗日以及许多其他人而达到全盛。平面和空间中坐标变换公式的线性特征——费马不可能没有已经察觉到这一点——例如由欧拉加以突出（{[}8(a){]}第二至第三章及第四章附录），他在这里为平面曲线和曲面按其次数进行分类奠定了基础（这一分类之所以不变，正是因为这些公式的线性性质）；也是他（同上引书，第十八章）引入了“仿射”一词，用以描述可以通过变换 \(x' = ax\) , \(y' = by\) 从其中一条导出另一条的曲线之间的关系（但在这个定义中，他没有看出任何几何上不变的东西，因为这个定义仍然依赖于坐标轴的特定选择）。稍后我们看到拉格朗日 {[}9(a){]} 用了整整一篇回忆录来专门讨论三维解析几何中典型的线性和多重线性问题，这篇回忆录长期享有应得的盛名。大约在这个时期，与“求经过给定点的平面曲线”这一线性问题相关，行列式的概念开始成形，最初在某种程度上是经验性的，出现在克拉默 {[}10{]} 和贝祖 {[}11{]}的著作中；随后这一概念被多位作者发展，其本质性质最终由柯西 {[}13{]} 和雅可比 {[}16(a){]} 确定下来。

另一方面，虽然数学家们对一次方程略有轻视，但微分方程的求解却被视为一个头等重要的问题；因此很自然，在这些方程中，线性方程（无论系数是否为常数）很早就被区分出来，而对其研究也有助于突出线性及其相关性质。这一点在拉格朗日 \([9(b)]\) 和欧拉 \([8(b)]\) 的工作中肯定可以看到，至少就齐次方程而言是如此；因为在这些作者看来，没有必要指出非齐次方程的通解是特解与相应齐次方程的通解之和，他们也未利用这一原理（尽管达朗贝尔已知此原理）；我们在此还要指出，当他们陈述n阶齐次线性方程的通解是n个特解的线性组合时，他们并没有补充说这些特解必须是线性无关的，也没有努力使后一概念明确化；似乎只有柯西在综合工科学校的教学才对这些点以及其他许多点有所阐明（{[}14{]}，第573—574页）。但拉格朗日（同上引书）也已经引入了 \(L^{*}(y)=0\) ，即线性微分方程 \(L(y)=0\) 的伴随方程（诚然纯粹是通过计算引入的，而且没有给它命名），这是对偶性的一个典型例子，其依据是关系式

\[
\int z \mathrm{L} (y) d x = \int \mathrm{L} ^ {*} (z) y d x,
\]

该关系式在 y 和 z 于积分区间端点处为零时成立；更确切地说，在高斯明确给出3个变量中线性替换的转置之前30年，我们在这里看到了毫无疑问的第一个“函数算子”L*的例子，即由双线性函数（这里是积分 \(\int yz dx\) ）给出的算子 L 的转置或“伴随”。

与此同时，同样是在拉格朗日 \([9(c)]\) 那里，线性替换（起初是2个和3个变量的）正在征服算术的过程中。显然，当x和y取所有整数值时，函数 \(F(x,y)\) 的值集在对其施行行列式为1、系数为整数的线性替换时不会改变；拉格朗日正是以这一基本观察为基础，建立了数用型表示的理论以及型的化简理论；而高斯则通过一个我们今天已难以体会其大胆程度的步骤，分离出等价的概念和型的类概念（参见I的历史注记）；在这个问题上，他认识到某些与线性替换有关的基本原理的必要性，并特别引入了转置或伴随的概念（{[}12(a){]}，第304页）。从这一刻起，对二次型（先是2个、3个变量，后来是n个变量）的算术研究和代数研究，对与之密切相关的双线性型的研究，以及更近期将这些概念推广到无穷多个变量的研究，直到今天一直是线性代数进步的最丰富源泉之一（参见IX的历史注记）。

但也许一个更具决定性的进步是高斯在同一部《算术研究》（参见I的历史注记）中创立了有限交换群理论，该理论在那里以四种不同的方式出现：在模m的整数加法群中（m为整数）、在模m的与m互素的整数乘法群中、在二元二次型的类群中，最后在m次单位根的乘法群中；而且，正如我们已经指出的，高斯显然是把所有这些群作为交换群，或者更确切地说，作为Z上的模来处理的，并研究它们的结构、同构关系等等。在“复整数” \(a + bi\) 所构成的模中，我们后来看到他研究一个Z上的无穷模，他无疑已经察觉到这个模与椭圆函数的周期模（由他在复域中发现）之间的同构；无论如何，这一思想已经清晰地出现在雅可比的工作中，例如在他关于具有3个周期的函数不可能性的著名证明以及他对阿贝尔积分反演问题的看法中 {[}16(b){]}，不久便导致了克罗内克的一系列定理（参见《一般拓扑学》历史注记，第七章）。

在这里，另一条支流汇入了我们一直试图追溯其流向与偶尔蜿蜒的河流，而这条支流长期以来一直潜藏于地下。正如稍后将在更详细的内容中阐述的那样（第九章的历史注记），上个世纪所理解的“纯粹”几何，即本质上不使用坐标的平面与空间射影几何，是由德扎格 {[}6{]} 在十七世纪创立的，其思想的真正价值为费马所认识，并为帕斯卡付诸实践，但随后即被埋没于遗忘之中，被解析几何的辉煌进展所遮蔽；直到十八世纪末，随着蒙日，继而庞斯莱及其对手布里安雄和沙勒的出现，它才得以复兴，有时完全且自觉地与解析方法分离，有时（尤其在德国）则与之紧密交织。如今，射影变换无论从何种观点（综合的或解析的）来看，当然都只是对射影坐标或“重心”坐标的线性替换；圆锥曲线理论（在十七世纪）以及后来的二次曲面理论——这一学派在很长一段时间内主要关注其射影性质——都只是二次型理论，而二次型与线性代数的密切联系我们此前已经指出。在这些概念之上又增添了配极概念：同样由德扎格所创，极与极线理论在蒙日及其后继者手中，并很快以对偶原理之名，成为变换几何定理的有力工具；如果说不能断言在那段时期人们已察觉到它与伴随微分方程的关系（这些关系在世纪末由平凯莱指出），那么至少沙勒没有忽略 {[}17{]} 它与互余球面三角形概念之间的联系——这一概念早在十六世纪就由维埃特（{[}3{]}，第428页）和斯涅尔引入球面三角学。但射影几何中的对偶性只是向量空间对偶性的一个方面，还需考虑从仿射空间过渡到射影空间（射影空间是仿射空间在“标量乘法”关系下的商空间）时所施加的修正。

十九世纪，在我们历史上比任何时期都更涌现出众多一流的数学家；即便只限于最显著的特征，要在几页之内描述这些思想运动汇聚于他们手中所产生的全部成果，也是困难的。一方面有纯粹综合的方法，那是一种正统拥护者自陷折磨的“普罗克鲁斯忒斯之床”；另一方面有与任意强加于空间的坐标系相关的分析方法；在这两者之间，人们很快便感到需要一种几何演算，这是莱布尼茨梦想却未创立、卡诺仅作了不完善勾画的演算：首先出现的是向量的加法，隐含于高斯关于虚数的几何表示及其对初等几何的应用之中（参见《一般拓扑学》历史注记，第八章），由贝拉维蒂斯以“等值法”之名加以发展，并在格拉斯曼、莫比乌斯和哈密顿那里获得最终形式；与此同时，莫比乌斯以“重心演算”之名给出了一个适合射影几何需要的版本 {[}18{]}。

在同一时期，由同一批人，从平面和“普通”空间迈向n维空间的这一步——一旦踏上这条道路便显得如此自然，且已由费马预示——被迈出了；这确实是不可避免的一步，因为能在两个或三个变量中用几何解释的代数现象，对于任意数量的变量依然成立；因此，在使用几何语言时，若将自己限制在二维或三维，对现代数学家而言，将如同那始终阻碍希腊人把数的概念扩展到不可公度量的比值的桎梏一样令人厌烦。于是，关于n维空间的语言和观念几乎同时在各方出现，在高斯的工作中尚显隐晦，在下一代数学家的著作中则清晰明了；而他们运用这些概念时信心的大小，或许与其说取决于他们的数学倾向，不如说取决于他们的哲学乃至纯粹实际的眼光。无论如何，凯莱和格拉斯曼大约在1846年已极其自如地处理这些概念（而且凯莱说——这与格拉斯曼恰恰相反（{[}22(a){]}，第321页）——是“不借助任何形而上学概念”）；凯莱从未远离解析解释和坐标，而格拉斯曼从一开始的工作中，n维空间中向量的加法及几何方面就占据了上风，最终导致了我们稍后将讨论的那些发展。

与此同时，高斯所给予的推动力正以两种不同的方式引导数学家们走向对代数或“超复数系”的研究。一方面，人们不可避免地试图以不同于引入“虚数单位” \(i = \sqrt{-1}\) 的方式来扩展实数域，或许由此开辟出与复数域同样丰饶的更为广阔的领域。高斯本人确信（{[}12(b){]}，第178页）只要人们想要保留复数的主要性质——用现代语言来说，即那些使它成为交换域的性质——这样的扩展就是不可能的；并且，无论受其影响还是独立地，他的同时代人似乎也持有这一信念，而这一信念直到很久以后才由魏尔斯特拉斯 {[}23{]} 以一条精确的定理加以证明。但是，一旦复数的乘法被解释为平面上的旋转，那么，如果提议将这一思想推广到三维空间（由于空间中的旋转构成非阿贝尔群），就必须考虑非交换的乘法；这正是哈密顿†发现四元数 {[}20{]} 时的指导性思想之一；四元数是非交换域的第一个例子。这个例子的独特性（正如后来弗罗贝尼乌斯所证明的，它是唯一能在实数域上构造的例子）在一定程度上限制了它的重要性，尽管——或者也许正是因为——围绕它形成了一个狂热的“四元数主义者”学派：这是一种奇怪的现象，后来在格拉斯曼的工作周围再次出现，随后又被那些从哈密顿和格拉斯曼那里提炼出所谓“向量分析”的通俗化者们所重演。稍后，格雷夫斯和凯莱放弃了结合性，构造了“凯莱数”，但这并未开辟出什么特别有趣的路径。但在西尔维斯特引入矩阵并（虽未命名但）清楚地定义了它们的秩 {[}21{]} 之后，又是凯莱 {[}22(b){]} 创立了矩阵的演算，并且并非没有注意到（一个后来常常被忽视的基本事实）矩阵只是线性替换的缩写记号，正如高斯将二次型 \(aX^{2} + 2bXY + cY^{2}\) 用记号 \((a, b, c)\) 来表示一样。这只是西尔维斯特和凯莱在行列式及其相关领域丰富产出中的一个方面——当然对我们来说是最有趣的一个方面——这种产出充满了巧妙的恒等式和令人印象深刻的计算。

此外（在其他事情中）格拉斯曼发现了实数上的一个代数，即至今仍以他命名的外代数。他的工作甚至早于哈密顿的工作 \([19(a)]\) ，是在几乎完全孤立的学术环境中完成的，长期鲜为人知，无疑是因为其独创性，也因为其开始时笼罩着的哲学迷雾，例如最初使莫比乌斯望而却步。受与哈密顿类似但更具重要性的关切所驱动（并且他很快意识到这些关注与莱布尼茨的相同），格拉斯曼构建了一个庞大的代数-几何大厦，建立在对n维向量空间的几何的或“内蕴的”（已经或多或少公理化的）概念之上；在他得到的较为初等的结果中，我们可举出例如向量线性无关的定义、维数的定义以及基本关系式

\[
\dim \mathrm{V} + \dim \mathrm{W} = \dim (\mathrm{V} + \mathrm{W}) + \dim (\mathrm{V} \cap \mathrm{W})
\]

（同上引书，第209页；参见{[}19(b){]}，第21页）。但尤其是多向量的外乘法，然后是内乘法，为他提供了工具，使他得以轻松地首先处理线性代数本身的问题，然后处理与欧几里得结构相关的问题，即向量的正交性（在那里他找到了对偶性的等价物，而他本人并不拥有对偶性概念）。

高斯在超复数系统研究中开辟的另一条路径是从复整数 \(a + bi\) 出发的；在这些之后，自然地出现了环Z上的代数或更一般的超复数系统，以及有理数域Q上的代数或更一般的超复数系统，首先是由单位根生成的、高斯已经设想过的那类系统，然后是代数数域和代数整数模：前者是库默尔工作的主要课题，后者的研究由狄利克雷、埃尔米特、克罗内克和戴德金着手进行。在这里，与实数上的代数情况相反，不必放弃交换域的任何特征性质，并且在整个十九世纪，人们的注意力都局限于后者。但线性性质以及例如寻找域的整数基（这是判别式的一般定义所必需的）在许多地方起着至关重要的作用；并且至少在戴德金那里，这些方法注定要变得典型地“超复数化”；此外，戴德金本人虽然没有给自己提出一般代数的问题，但他意识到其工作的这一特征，以及它们与例如魏尔斯特拉斯关于实数上超复数系统的结果之间的联系（{[}24{]}，特别是第2卷第1页）。与此同时，代数数域中乘法单位群结构的确定——由狄利克雷在一些著名的笔记 {[}15{]} 中、并且几乎同时由埃尔米特完成——对于澄清关于Z上模、其生成系及其基（当这样的基存在时）的观念至关重要。然后，戴德金在代数数域中定义了理想的概念（作为该域整数环上的模），而克罗内克在多项式环中以“模系统”为名引入了一个等价的概念，这给出了比Z更一般的环上模的最初例子；并且在同一些作者以及后来希尔伯特的工作中，在特殊情况下，带算子的群的概念被缓慢地分离出来，并且总是有可能从这样的群构造出适当定义的环上的模。

与此同时，对二次双线性形式及其“化简”（或等价地，对矩阵及其“不变量”）的算术-代数研究，导致了求解线性方程组的一般原理的发现，这些原理由于缺乏秩的概念而未被雅可比所察觉。† 关于整数系数线性方程组的整数解问题，先由埃尔米特在特殊情形下着手解决，随后由H. J. 史密斯 {[}25{]} 在完全一般性下加以解决；后者的结果直到1878年才由弗罗贝尼乌斯重新发现，当时正值克罗内克发起并由魏尔斯特拉斯也参与其中的一项庞大研究计划；顺便提一下，在这些工作过程中，克罗内克为实（或复）系数线性方程组的定理给出了最终形式，而这些定理也由《爱丽丝漫游奇境记》的著名作者以他特有的细致入微的方式在一本晦涩的手册中加以阐明；至于克罗内克本人，则不屑于发表这些结果，而将其留给同事和弟子们；就连“秩”这个词本身也是由弗罗贝尼乌斯引入的。同样在柏林大学的教学过程中，克罗内克 {[}26{]} 和魏尔斯特拉斯引入了行列式的“公理化”定义（作为n维空间中n个向量的交错多重线性函数，并通过规范化使其在单位矩阵处取值为1），该定义等价于由格拉斯曼演算导出的定义，也等价于本专著所采用的定义；同样在他的课程中，克罗内克在未感到需要为其命名的同时，以仍非内蕴的形式引入了空间的张量积和矩阵的“克罗内克”积（即由给定的线性替换分别作用于各因子而在张量积上诱导出的线性替换）。

这项研究离不开凯莱、埃尔米特和西尔维斯特所创立的“不变量理论”（即埃尔米特后来在书信中提到的“不变量三杰”），而从现代观点来看，这一理论首先是线性群的表示理论。在此，作为射影几何中对偶性的代数等价物，出现了共变变量组与反变变量组之间的区分，即空间中的向量与对偶空间中的向量；并且，在注意力首先转向低次形式、随后转向任意次形式（涉及2个和3个变量）之后，几乎紧接着就考察了多组“共变”或“反变”变量下的双线性形式，进而是多线性形式，这等价于引入张量；后者变得明确并得到普及，是在不变量理论的启发下，里奇和列维-奇维塔于1900年将“张量分析” {[}28{]} 引入微分几何之时；这一理论后来因“相对论”物理学家的使用而大为盛行。不变量理论、微分几何与偏微分方程理论（尤其是所谓的普法夫问题及其推广）之间日益深入的相互交融，再次缓慢地引导几何学家考虑微分的交错双线性形式，特别是1次形式的“双线性协变”（1870年由利普希茨引入，随后由弗罗贝尼乌斯研究），最终促成了E. 嘉当 {[}29{]} 与庞加莱 {[}30{]} 对外微分形式演算的创立。庞加莱引入后者，是为了构造他的积分不变量，将其作为多重积分中出现的表达式；而嘉当，无疑受到其代数研究的指引，以更形式化的方式引入它们，但并未忽视其演算的代数部分与格拉斯曼的外积乘法完全一致（由此得名，他采用了这一名称），从而最终恢复了后者工作的应有地位。将外微分形式翻译成张量演算的记号，也立即显示出它们与反对称张量的联系，一旦采用纯代数的观点，这表明它们之于交错多重线性形式，正如协变张量之于任意多重线性形式；这一方面在现代线性群表示理论中得到了进一步阐明；因此，例如，魏尔斯特拉斯和克罗内克给出的行列式定义与格拉斯曼演算所得定义之间的实质同一性得到了确认。

我们由此进入现代时期，公理化方法和结构概念（起初只是模糊地感知，直到最近才被定义）使我们能够分离此前一直纠缠不清的概念，将模糊或留待直觉的内容加以明确表述，并以适当的普遍性证明那些仅在特殊情形下已知的定理。佩亚诺，公理化方法的创立者之一，也是最早充分赏识格拉斯曼工作的数学家之一，早在1888年（{[}27{]}，第九章）就给出了实数域上向量空间（有限维或无限维）的公理化定义，并以完全现代的记号给出了一个这样的空间到另一个空间的线性映射的定义；稍后，平凯莱试图将如此构想的线性代数应用于函数理论，诚然这一方向并未产生丰硕成果；但至少他的观点使他能够将“拉格朗日伴随”识别为线性映射转置的一个特例：这一点很快在希尔伯特及其学派关于希尔伯特空间及其在分析中应用的难忘工作中，对偏微分方程和常微分方程而言，都表现得更加清晰。正是在那个场合，托普利茨 {[}31{]}还引入了（但借助坐标）实数域上最一般的向量空间，并作出了一个基本观察：证明线性代数的主要定理并不需要行列式理论，这使得这些定理可以毫无困难地推广到无穷维空间；他还指出，如此理解的线性代数当然可以应用于任意交换基域。

另一方面，随着巴拿赫于1922年引入以他命名的空间 \(^{\dagger}\) ，人们遇到了与其对偶空间不同构的空间，尽管这出现在一个既是拓扑的又是代数的问题中。即使在有限维向量空间与其对偶空间之间，也不存在由结构决定的“典范”同构，这一点早已反映在共变与反变的区分中。然而似乎毫无疑问，空间与其对偶空间之间的区分只是在巴拿赫及其学派的工作之后才被最终确立；同样在这些工作中，余维数概念的重要性也显现出来。至于一个空间的向量子空间与其对偶空间的向量子空间之间的对偶性或“正交性”，今天表述它的方式不仅与伽罗瓦理论主要定理的现代表述（参见《代数》V）以及局部紧阿贝尔群中所谓的庞特里亚金对偶性有着表面的相似；后者可追溯到韦伯，他在算术研究过程中于1886年为有限群奠定了其基础；在伽罗瓦理论中，子群与子域之间的“对偶性”在戴德金和希尔伯特的工作中成形；而向量子空间之间的正交性则明显地首先来源于射影几何中线性簇之间的对偶性，也来源于欧几里得空间或希尔伯特空间中完全正交簇的概念及其性质（其名称即由此而来）。所有这些线索在当代时期被重新汇聚，在诸如
E. 诺特、阿廷和哈塞等代数学家以及庞特里亚金和惠特尼等拓扑学家手中（并非没有相互影响），从而在每一个这些领域中达到了本专著所阐述的知识阶段。

与此同时，进行了一种批判性的审视，其目的在于在每一点上消除那些并非完全不可或缺的假设，尤其是那些会阻碍某些应用的假设。因此，人们认识到在向量空间的概念中用环代替域的可能性，并通过创造模的一般概念，同时处理这些空间、阿贝尔群、克罗内克、魏尔斯特拉斯、戴德金和施泰尼茨已经研究过的特殊模，甚至带算子的群，例如将若尔当-赫尔德定理应用于它们；同时，随着右模与左模的区分，我们进入了非交换的情形，这源于美国学派（韦德伯恩、迪克森）尤其是德国学派（E. 诺特、E. 阿廷）对代数理论的现代发展。

